% ---------- ¶4  the assumption audit ----------
The unraveling prediction rests on four assumptions: (a)~buyers know sellers possess private
information about quality, (b)~disclosure is costless, (c)~disclosure is credible, and
(d)~buyers are homogeneous in how much they value the disclosed attribute. Three of the four are
either directly satisfied in this setting or can be dispensed with. Assumption~(c) holds by
construction after July~2021, since the DPE is issued by a certified third-party diagnostician,
archived in the registry of the \emph{Agence de la Transition \'Ecologique} (ADEME), and legally
binding on the seller. Assumption~(b) is close to satisfied, because a DPE is a one-off outlay of
roughly 100--250~\euro{} according to ADEME's own estimates, negligible relative to
residential transaction values, and it is the seller's obligation rather than the buyer's.
Assumption~(a), which the theoretical literature treats as a maintained hypothesis, we do not
need to assume at all. By matching listings to the ADEME registry we observe directly which
sellers hold a certificate at the moment of listing, and among the apartment listings we can
match at the address level, 48.8\% hold none.
Silence in this market is therefore not always a choice, a possibility first identified by
\citet{dye_disclosure_1985} and formalised by \citet{jung_kwon_1988}, and here it is measured
rather than inferred.
